% =============================================================================
%  Theory of Generated Space (TGP): Core
%  A self-contained preprint defining axioms, substrate, effective field,
%  principal equations, and the status of theorems and open problems.
%
%  Mateusz Serafin,  April 2026
%  To be published on Zenodo.
% =============================================================================
\documentclass[11pt,a4paper]{article}

\usepackage[utf8]{inputenc}
\usepackage[T1]{fontenc}
\usepackage[english]{babel}
\usepackage{amsmath,amssymb,amsfonts,amsthm}
\usepackage{mathtools}
\usepackage[hidelinks]{hyperref}
\usepackage{geometry}
\usepackage{longtable,booktabs,tabularx,array}
\usepackage[shortlabels]{enumitem}
\usepackage{xcolor}
\usepackage[protrusion=true,expansion=false]{microtype}
\geometry{margin=2.5cm}

% --- TGP symbols -------------------------------------------------------------
\renewcommand{\Phi}{\varPhi}
\newcommand{\PhiZero}{\varPhi_{0}}
\newcommand{\Nzero}{N_{0}}
\newcommand{\Ffield}{\mathcal{F}}
\newcommand{\lP}{\ell_{\mathrm{P}}}
\newcommand{\Zp}{\mathbb{Z}_{2}}
\newcommand{\DD}{\mathcal{D}}
\newcommand{\NN}{\mathcal{N}}

% --- Theorem environments ----------------------------------------------------
\newtheorem{axiom}{Axiom}
\newtheorem{definition}[axiom]{Definition}
\newtheorem{theorem}[axiom]{Theorem}
\newtheorem{proposition}[axiom]{Proposition}
\newtheorem{lemma}[axiom]{Lemma}
\newtheorem{corollary}[axiom]{Corollary}
\newtheorem{remark}[axiom]{Remark}
\newtheorem{hypothesis}[axiom]{Hypothesis}

\title{%
  \textbf{Theory of Generated Space (TGP):}\\[4pt]
  \large A minimal core of axioms, substrate, and effective field%
}
\author{Mateusz Serafin\\[2pt]
\small\textit{Independent researcher}\\[-2pt]
\small\texttt{environmenstefan@gmail.com}}
\date{April 2026}

\begin{document}
\maketitle

\begin{abstract}
We present the core of the \emph{Theory of Generated Space} (TGP):
a scalar-field theory in which space is not a background arena but is
\emph{constituted} by mass-energy sources through a shared field $\Phi$.
The theory is built on four axioms
(absolute reference state $\Nzero$, $\Zp$-symmetric relational substrate,
matter as source, instability of $\Nzero$) and a single dynamical degree of freedom.
From these inputs we derive:
(i) a unique second-order kinetic operator with coefficient $\alpha=2$,
as a classification theorem within axioms (C1)--(C3) of
Sec.~\ref{sec:field} (its derivation from the discrete substrate
$H_{\Gamma}$ is stated as an open problem, OP-6);
(ii) an effective spatial metric of exponential form $g_{ij}=e^{+2U/c_{0}^{2}}\delta_{ij}$
that reproduces all post-Newtonian parameters of general relativity exactly
($\gamma_{\mathrm{PPN}}=\beta_{\mathrm{PPN}}=1$);
(iii) three interaction regimes (long-range attraction, intermediate repulsion,
short-range confinement) from a single non-linear potential $U(\Phi)$;
(iv) identically luminal gravitational-wave propagation
$c_{\mathrm{GW}}=c_{0}$, consistent with GW170817.
TGP is \emph{not} a scalar-tensor theory: the metric is algebraically determined
by $\Phi$, and the action contains no Einstein--Hilbert term.
This document defines the minimal, self-contained core; applications
(superconductivity, lepton masses, black-hole shadows, dark energy) are
developed in companion papers that cite this reference.
We state twenty-one proven or numerically verified results and ten
explicit open problems, so that the theory is falsifiable and extendable
on a well-defined basis.
\end{abstract}

\noindent\textbf{Keywords:}
generated space, emergent gravity, scalar field, $\Zp$ substrate,
PPN parameters, effective metric.

\tableofcontents
\newpage

% =============================================================================
\section{Introduction}\label{sec:intro}
% =============================================================================

The standard model of physics, combined with $\Lambda$CDM cosmology and
general relativity, is phenomenologically successful but treats space as
a pre-existing manifold on which fields propagate. Approximately thirty-five
free parameters govern the combined description, and the origin of the
small cosmological constant, the hierarchy of lepton masses, and the
mechanism of high-temperature superconductivity remain open.

We propose an alternative foundation in which space is not pre-given.
Instead, every source of mass-energy \emph{produces} a local contribution
to a shared scalar field $\Phi$, whose value measures the density of generated
space. The metric, causality and Newtonian gravity emerge as effective
descriptions of the configuration of $\Phi$; the Planck length is the only
invariant carried across scales.

The purpose of this paper is to \emph{define} the theory --- its axioms,
substrate, effective field, and canonical equations --- in a minimal,
self-contained form, and to delineate explicitly what has been proven and
what remains open. We consciously avoid a survey of applications: these are
developed in companion papers that cite this reference (see
Sec.~\ref{sec:applications}).

\paragraph{What TGP is not.} TGP is not a scalar-tensor theory: the
effective metric is algebraically determined by $\Phi$, and the action
contains no Einstein--Hilbert term. TGP is not a programme in which space
emerges from information or entanglement (AdS/CFT, ER\,=\,EPR): here space
emerges from \emph{matter}. TGP is not an extension of general relativity
with additional fields; it is a reformulation in which GR is recovered in
a well-defined regime.

\paragraph{What TGP is.} A single scalar degree of freedom $\Phi$,
coarse-grained from a $\Zp$-symmetric relational substrate, satisfying a
single nonlinear field equation sourced by matter, and inducing an
exponential effective metric. From this minimal structure flow all
predictions of the theory.

\paragraph{Scope of this paper.} Sections~\ref{sec:axioms}--\ref{sec:metric}
establish the axiomatic base, the substrate-to-field derivation, the
effective field theory, and the emergent metric. Section~\ref{sec:regimes}
exhibits the three physical regimes. Section~\ref{sec:status} enumerates the
status of theorems and open problems --- the central artifact of this paper
from a reproducibility standpoint. Section~\ref{sec:applications} is a brief
forward reference to companion papers.

\paragraph{Main open caveat.} The passage from the discrete substrate
Hamiltonian $H_{\Gamma}$ to the continuum TGP action is \emph{not yet}
a rigorous theorem. Numerical block-spin Monte Carlo and FRG LPA'
computations on a one-dimensional substrate chain support convergence
to the $\alpha=2$ fixed point (see row~21 of
Sec.~\ref{sec:status}), but a Banach/$\Gamma$-convergence proof remains
open (OP-6). Throughout this paper we therefore present the continuum
field theory of Sec.~\ref{sec:field} as a well-posed effective theory
whose numerical coarse-graining has been checked, while flagging the
rigorous continuum-limit statement as outstanding.

% =============================================================================
\section{Axioms}\label{sec:axioms}
% =============================================================================

TGP rests on four axioms. Axioms \ref{ax:N0}--\ref{ax:P3} concern the
reference state and its instability; Axiom~\ref{ax:source} states how
matter couples to space. A fifth, derived statement (the $\Zp$ symmetry
of the substrate) appears as a proposition rather than a postulate.

\begin{axiom}[Absolute reference state]\label{ax:N0}
There exists a reference state $\Nzero$ -- \emph{absolute nothingness} --
characterized by the simultaneous vanishing of the state function
$\Ffield$ and the field $\Phi$:
\begin{equation}\label{eq:N0}
  \Nzero := \bigl\{\emptyset,\;\Ffield\bigr\},
  \qquad
  \Ffield\big|_{\Nzero}=0,
  \qquad
  \Phi\big|_{\Nzero}=0.
\end{equation}
In $\Nzero$ there is no metric, no matter, no time, no energy.
$\Nzero$ is \emph{not} the quantum vacuum and \emph{not} empty space;
it is the absence of every physical metric structure.
\end{axiom}

\begin{axiom}[Relational substrate]\label{ax:substrate}
Underlying both the $\Nzero$ state and all physical configurations is a
fundamental, non-emergent relational substrate
\begin{equation}\label{eq:Gamma}
  \Gamma=(V,E),
\end{equation}
with vertices $V$ and edges $E$, possessing adjacency relations and
connectivity but \emph{no a priori metric}. Each node $i\in V$ carries an
amplitude operator $\hat{s}_{i}$. The substrate Hamiltonian is
\begin{equation}\label{eq:H-Gamma}
  H_{\Gamma}=\sum_{i}\!\Bigl[
    \tfrac{\hat{\pi}_{i}^{2}}{2\mu}
    +\tfrac{m_{0}^{2}}{2}\hat{s}_{i}^{2}
    +\tfrac{\lambda_{0}}{4}\hat{s}_{i}^{4}
  \Bigr]
  - J\!\!\sum_{\langle ij\rangle}\! A_{ij}\,\hat{s}_{i}\hat{s}_{j},
\end{equation}
with $J>0$, $\lambda_{0}>0$, and $A_{ij}\in\{0,1\}$ the adjacency matrix.
\end{axiom}

\begin{proposition}[$\Zp$ symmetry]\label{prop:Z2}
The Hamiltonian~\eqref{eq:H-Gamma} is invariant under the chiral inversion
$\hat{s}_{i}\mapsto-\hat{s}_{i}$ for all $i\in V$. This discrete $\Zp$
symmetry is the fundamental symmetry of TGP; the state function $\Ffield$
measures the extent of its spontaneous breaking.
\end{proposition}

\begin{axiom}[Matter generates space]\label{ax:source}
Every source of mass-energy $\rho_{a}$ contributes to the field $\Phi$
according to
\begin{equation}\label{eq:source}
  \DD[\Phi_{a}]\;=\;q_{a}\,\rho_{a},
\end{equation}
where $q_{a}>0$ is the \emph{spatial generation constant} of particle
type $a$ and $\DD$ is the covariant operator governing the dynamics of
$\Phi$ (made explicit in Sec.~\ref{sec:field}).
\end{axiom}

\begin{axiom}[Instability of $\Nzero$]\label{ax:P3}
The state $\Ffield=0$ ($\Nzero$) is unstable: its measure in configuration
space vanishes,
\begin{equation}
  \mu\bigl(\{\sigma\in\Omega : \Ffield[\sigma]=0\}\bigr)=0.
\end{equation}
Any departure from exact $\Zp$ symmetry leads locally to $\Ffield>0$ and
hence to simultaneous matter and generated space.
\end{axiom}

\begin{proposition}[Instability from substrate thermodynamics]\label{prop:instability}
Let $\mu_{\beta}(\sigma)\propto e^{-\beta H_{\Gamma}[\sigma]}$ be the Gibbs
measure for~\eqref{eq:H-Gamma}. In the thermodynamic limit $|V|\to\infty$
there exists $\beta_{c}>0$ such that
$\mu_{\beta}\bigl(\{\Ffield=0\}\bigr)=0$ for all $\beta>\beta_{c}$; the
physical substrate ($\beta\to\infty$) is always in the spontaneously
broken regime. Thus Axiom~\ref{ax:P3} is a \textbf{consequence} of
Axioms~\ref{ax:substrate}--\ref{ax:source} together with standard Ising
thermodynamics, not an independent postulate.
\end{proposition}

\begin{remark}[Three-level ontology]
These axioms generate a three-level hierarchy:
(0) the relational substrate $\Gamma$ is fundamental and non-emergent;
(1) the field $\Phi$ is a coarse-grained \emph{metric-phase parameter} of
the substrate, with $\Phi=0$ the non-metric phase and $\Phi>0$ the metric
phase;
(2) the effective geometry $g_{\mu\nu}$ arises from coarse-graining
level~(1). The transition $\Phi=0\to\Phi>0$ does not create space from
nothing: it switches on the metric response of an already existing
substrate, in the same sense that freezing water does not create
molecules but changes their phase.
\end{remark}

% =============================================================================
\section{From substrate to field}\label{sec:substrate-to-field}
% =============================================================================

The relation between the microscopic substrate $\Gamma$ and the macroscopic
field $\Phi$ is established by coarse-graining. Given a block
$\mathcal{B}(\mathbf{x})$ of lattice sites centred on $\mathbf{x}$,
\begin{equation}\label{eq:Phi-def}
  \Phi(\mathbf{x})
  \;\propto\;
  \langle\hat{s}^{2}\rangle_{\mathcal{B}(\mathbf{x})}.
\end{equation}
In the mean-field approximation one has the identity
\begin{equation}\label{eq:F-s-relation}
  \Ffield_{i}[\Omega]
  \;=\;\langle\Omega|\,\hat{s}_{i}\,|\Omega\rangle^{2},
\end{equation}
so that $\Ffield=0$ iff $\langle\hat{s}_{i}\rangle=0$ (unbroken $\Zp$,
i.e.\ $\Nzero$), while $\Ffield>0$ signals spontaneous breaking and the
simultaneous appearance of matter and space. A Migdal--Kadanoff
coarse-graining applied to~\eqref{eq:H-Gamma} produces the continuum field
equation for $\Phi$ that we state and analyse in Sec.~\ref{sec:field}.

\begin{remark}[One substrate, two projections]
Coarse-graining $H_{\Gamma}$ yields \emph{two} emergent fields from a
single microscopic object:
\begin{enumerate}[(i),nosep]
  \item the scalar metric-phase parameter
    $\Phi(\mathbf{x})\propto\langle\hat{s}^{2}\rangle_{\mathcal{B}}$,
    which governs the entire gravitational/cosmological sector;
  \item a spatial-correlation tensor
    $\sigma_{ab}(\mathbf{x})\propto
    \langle\hat{s}_{i}\,\hat{s}_{i+\hat{e}_{a}}\rangle_{\mathcal{B}}^{\mathrm{TF}}$,
    carrying spin-2 and underlying tensor gravitational-wave modes.
\end{enumerate}
Both follow from the \emph{same} $H_{\Gamma}$ and are not independent
postulates. This is crucial: TGP is not a scalar-tensor theory, because
$\Phi$ and $\sigma_{ab}$ are two projections of one substrate rather than
two independent fields.
\end{remark}

% =============================================================================
\section{Effective field theory}\label{sec:field}
% =============================================================================

We use the dimensionless variable $\varphi \equiv \Phi/\PhiZero$.

\subsection{Uniqueness of the kinetic operator}

\begin{theorem}[Uniqueness of $\DD_{\mathrm{kin}}$ within class (C1)--(C3), $\alpha=2$]\label{thm:alpha2}
Among all local second-order scalar kinetic operators of the form
\begin{equation}
  \DD[\varphi]
  = K(\varphi)\,\nabla^{2}\varphi
  + \tfrac{1}{2}K'(\varphi)(\nabla\varphi)^{2},
\end{equation}
imposing
\begin{enumerate}[(C1),nosep]
  \item $K(\varphi)$ smooth and depending only on $\varphi$ (locality);
  \item $\DD[\varphi]\to\nabla^{2}\varphi$ for $\varphi=\mathrm{const}$
    (Laplacian limit);
  \item covariance under conformal rescaling $\varphi\mapsto\lambda\varphi$
    of the substrate metric $h_{ij}=\varphi^{4}\,\delta_{ij}$
    ($\Phi$-covariance),
\end{enumerate}
selects uniquely $K(\varphi)=K_{\mathrm{geo}}\varphi^{4}$ with the
ratio $\alpha(\varphi)\equiv K'(\varphi)\varphi/[2K(\varphi)]=2$ (constant).
The kinetic operator is
\begin{equation}\label{eq:Dkin}
  \boxed{\;
    \DD_{\mathrm{kin}}[\varphi]
    \;=\;\nabla^{2}\varphi+\frac{2(\nabla\varphi)^{2}}{\varphi}
    \;=\;\frac{1}{3\varphi^{2}}\,\nabla^{2}(\varphi^{3}).
  \;}
\end{equation}
\end{theorem}

\begin{proof}[Proof sketch]
Condition (C2) fixes $K(1)=1$. The Euler--Lagrange equation of
$\mathcal{L}_{\mathrm{kin}}=-\tfrac{1}{2}K(\varphi)(\nabla\varphi)^{2}$
yields the ratio $\alpha(\varphi)=K'(\varphi)\varphi/[2K(\varphi)]$.
(C1) forces $\alpha$ to be a constant; (C3) fixes $\alpha=2$ through the
explicit computation of the Laplace--Beltrami operator on
$h_{ij}=\varphi^{4}\,\delta_{ij}$:
$\Delta_{h}\varphi=\varphi^{-4}\,\DD_{\mathrm{kin}}[\varphi]$ holds
iff $\alpha=2$. Hence $K=\varphi^{4}$ (up to the overall constant
$K_{\mathrm{geo}}$, conventionally $1$).
\end{proof}

\begin{remark}[What $\alpha=2$ means]
The exponent $\alpha=2$ is \emph{not} fitted to data: it is fixed by the
geometric coupling $K_{ij}=J(\varphi_{i}\varphi_{j})^{2}$ between nearest
neighbours on the substrate, under the three conditions (C1)--(C3). The
statement is therefore a uniqueness result \emph{within} the stated
class of local, $\Phi$-covariant second-order scalar operators; relaxing
any of (C1)--(C3) (non-locality, higher derivatives, or departures from
the conformal rescaling of the substrate metric) would open additional
channels that are not considered by TGP. The derivation of (C1)--(C3)
as emergent properties of the block-averaged substrate $\varphi\mapsto
\Phi=\varphi^{2}$ is itself still open (OP-6); one-loop LPA' flow shows
that the kinetic structure $K(\rho)\!\propto\!\rho$ is stable in the
continuum limit (row~16 of Sec.~\ref{sec:status}).
\end{remark}

\subsection{Action, potential, and field equation}

The complete TGP action is
\begin{equation}\label{eq:S-TGP}
  S_{\mathrm{TGP}}[\Phi,\psi_{m}]
  =\int d^{4}x\,\sqrt{-g_{\mathrm{eff}}}\,
    \Bigl[\mathcal{L}_{\mathrm{field}}(\Phi)
    +\mathcal{L}_{\mathrm{mat}}(\Phi,\psi_{m})\Bigr],
\end{equation}
with field Lagrangian
\begin{equation}\label{eq:L-field}
  \mathcal{L}_{\mathrm{field}}
  =-\tfrac{1}{2}\,K(\varphi)\,
    g_{\mathrm{eff}}^{\mu\nu}\partial_{\mu}\varphi\,\partial_{\nu}\varphi
  - U(\varphi),
  \qquad K(\varphi)=\varphi^{4},
\end{equation}
self-interference potential
\begin{equation}\label{eq:U-potential}
  U(\varphi)
  =\frac{\beta}{3}\varphi^{3}-\frac{\gamma}{4}\varphi^{4},
  \qquad \beta=\gamma
  \;\;\text{(vacuum condition $U'(1)=0$)},
\end{equation}
and minimal matter coupling
$\mathcal{L}_{\mathrm{mat}}=-\,q\,\varphi\rho/\PhiZero$. The volume element
is $\sqrt{-g_{\mathrm{eff}}}=c_{0}\varphi$ (see Sec.~\ref{sec:metric}).

\begin{theorem}[Full TGP field equation]\label{thm:field-eq}
Variation $\delta S_{\mathrm{TGP}}/\delta\Phi=0$ of~\eqref{eq:S-TGP} yields
\begin{equation}\label{eq:field-eq}
  \boxed{\;
    \nabla^{2}\Phi
    +\frac{2(\nabla\Phi)^{2}}{\Phi}
    +\beta\,\frac{\Phi^{2}}{\PhiZero}
    -\gamma\,\frac{\Phi^{3}}{\PhiZero^{2}}
    \;=\;-\,q\,\PhiZero\,\rho.
  \;}
\end{equation}
\end{theorem}

\begin{remark}[TGP is not a scalar-tensor theory]\label{rem:not-ST}
The action~\eqref{eq:S-TGP} contains \emph{no} Einstein--Hilbert term $R$.
The effective metric $g_{\mu\nu}^{\mathrm{eff}}$ is defined algebraically
by $\Phi$ (not an independent dynamical variable). The Einstein equations
arise as a geometric consistency condition of $g_{\mu\nu}^{\mathrm{eff}}$,
propagated by the Bianchi identity once satisfied at an initial time; they
are not postulated. The only dynamical degree of freedom is $\Phi$.
\end{remark}

\begin{remark}[Formal contrast with scalar-tensor theories]\label{rem:ST-formal}
Structurally TGP differs from the Brans--Dicke / Horndeski family on three
counts, each visible at the level of the action and the equations of
motion:
\begin{enumerate}[(i),nosep]
  \item \emph{Independent tensor d.o.f.\ absent.} In scalar-tensor
    theories the action is $S=\int\sqrt{-g}\,[f(\phi)R/16\pi G+\ldots]$
    with $g_{\mu\nu}$ varied independently. Here the only field varied
    is $\Phi$; the metric is algebraically determined by
    $g_{\mu\nu}^{\mathrm{eff}}[\Phi]$ through Eq.~\eqref{eq:exp-metric}.
  \item \emph{No conformal frame ambiguity.} Scalar-tensor theories
    admit Jordan/Einstein frame reparametrisations. TGP has a single
    frame: the substrate metric $h_{ij}=\Phi^{4}\delta_{ij}$ is fixed by
    the $\Zp$ substrate structure; no conformal redefinition of the
    ``physical'' metric is available.
  \item \emph{Kinetic operator not a free function.} In scalar-tensor
    theories $f(\phi)$ is an input function to be fitted. Here
    $K(\varphi)=\varphi^{4}$ is fixed by
    Theorem~\ref{thm:alpha2} under (C1)--(C3), not chosen by hand.
\end{enumerate}
These three properties preclude TGP from being recast as a standard
scalar-tensor theory in any frame; the GR limit is recovered only as an
effective low-curvature description, not as a conformal transformation
of a two-metric theory.
\end{remark}

\subsection{Ghost-freedom}\label{ssec:ghost}

The substrate coupling $K_{\mathrm{sub}}(\varphi)=\varphi^{4}$ is
strictly positive for all $\varphi>0$, so the kinetic term
$\tfrac{1}{2}\varphi^{4}(\nabla\varphi)^{2}\geq 0$ is positive-semidefinite.
TGP is ghost-free at three levels simultaneously:

\begin{center}\small
\begin{tabular}{@{}lll@{}}
\toprule
\textbf{Sector} & \textbf{Kinetic positivity} & \textbf{Mechanism}\\
\midrule
Fundamental ($\varphi$) & $\varphi^{4}(\nabla\varphi)^{2}>0$ & geometric\\
Mukhanov--Sasaki perturbations & $Q_{s}=\varphi_{\mathrm{bg}}^{4}>0$ & linearized\\
Solitons (full lattice) & $K_{\mathrm{sub}}=\varphi^{4}>0$ & substrate\\
\bottomrule
\end{tabular}
\end{center}

A logarithmic truncation $K=1+4\ln\varphi$ that appears in certain
expansions exhibits an apparent zero at $\varphi^{\ast}=e^{-1/4}$; this
is an artifact of the Taylor expansion of the exponent,
$\varphi^{4}=e^{4\ln\varphi}=1+4\ln\varphi+\dots$, not a physical
singularity.

% =============================================================================
\section{Emergent metric}\label{sec:metric}
% =============================================================================

The effective metric of TGP is not postulated. It is derived from three
inputs:
(i) node density of the substrate $\Rightarrow$ $g_{ij}=\psi\,\delta_{ij}$
with $\psi\equiv\Phi/\PhiZero$;
(ii) information budget $\Rightarrow f\cdot h=1$ (antipodal condition);
(iii) PPN consistency $\Rightarrow$ exponential form.

\subsection{Spatial part: node-density ansatz}

A coarse-graining cell $\mathcal{B}(\mathbf{x})$ containing $N_{B}$
lattice nodes has physical volume proportional to $\Phi(\mathbf{x})$:
\begin{equation}\label{eq:g-spatial}
  g_{ij}=\frac{\Phi}{\PhiZero}\,\delta_{ij}=\psi\,\delta_{ij},
  \qquad \psi\equiv\Phi/\PhiZero.
\end{equation}

\subsection{Temporal part: information budget}

The substrate has a fixed information capacity $\mathcal{B}=N_{B}s_{0}$ per
cell. Partitioning between spatial capacity $h(\Phi)$ and temporal
capacity $f(\Phi)$, the multiplicative conservation law
\begin{equation}\label{eq:fh-1}
  f(\Phi)\cdot h(\Phi)=1
\end{equation}
gives $f=\PhiZero/\Phi=1/\psi$.

\subsection{Exponential metric from PPN consistency}

\begin{theorem}[Selection of $h(\Phi)=\Phi$ and exponential form,
within the TGP ansatz class]\label{thm:metric}
Within the ansatz class
\begin{equation}
  g_{tt}=-c_{0}^{2}/\psi,\qquad g_{ij}=h(\Phi)\,\delta_{ij},
  \qquad fh=1\ \text{(antipodal projection)},
\end{equation}
and in the weak-field limit $\Phi=\PhiZero(1+2U/c_{0}^{2})$ with $U$ the
Newtonian potential, the power-law choice $h=(\Phi/\PhiZero)^{p}$ forces
$\beta_{\mathrm{PPN}}=2p$; PPN consistency $\beta_{\mathrm{PPN}}=1$ then
selects $p=1$ uniquely. Continuing the same power-law structure to
\emph{all} orders in $U$ requires the exponential closure
\begin{equation}\label{eq:exp-metric}
  \boxed{\;
    ds^{2}
    =-c_{0}^{2}\,e^{-2U/c_{0}^{2}}\,dt^{2}
    +e^{+2U/c_{0}^{2}}\,\delta_{ij}\,dx^{i}dx^{j},
    \qquad
    \psi=e^{+2U/c_{0}^{2}}.
  \;}
\end{equation}
Both forms $f=1/\psi$ and $f=e^{-2U/c_{0}^{2}}$ satisfy $fh=1$ and agree
at $O(U)$; the exponential form extends the agreement to all orders.
\end{theorem}

\begin{remark}[Scope of the selection statement]\label{rem:metric-scope}
Theorem~\ref{thm:metric} is a selection statement \emph{within} the TGP
ansatz class $fh=1$ and power-law $h(\Phi)$. It is not a uniqueness
statement against every conceivable $(f,h)$; a parametrically distinct
metric structure (e.g.\ bi-metric or tensor-valued) would not be in
competition here. Five independent arguments support the choice $p=1$
(substrate-density interpretation, volume element, Cassini/LLR PPN,
soliton mass ratio, information budget); see
\texttt{research/metric\_ansatz/} for details.
\end{remark}

\begin{corollary}[PPN parameters exact, within the ansatz of Theorem~\ref{thm:metric}]\label{cor:ppn}
As an algebraic consequence of the exponential ansatz~\eqref{eq:exp-metric}
(selected within the TGP ansatz class; see
Remark~\ref{rem:metric-scope}), expanding the metric and comparing with
the standard PPN expansion, one obtains
\begin{equation}\label{eq:ppn}
  \gamma_{\mathrm{PPN}}=1,\qquad \beta_{\mathrm{PPN}}=1
  \qquad\text{(exact to all orders in $U$).}
\end{equation}
All ten PPN parameters equal their GR values. The prediction is
consistent with every current solar-system test:
Cassini Shapiro-delay tracking gives
$|\gamma_{\mathrm{PPN}}-1|\leq 2.3\times 10^{-5}$~\cite{BertottiCassini},
and Lunar Laser Ranging (Nordtvedt effect) gives
$|\beta_{\mathrm{PPN}}-1|\leq 6\times 10^{-5}$~\cite{WilliamsLLR};
a dedicated cross-check on all ten PPN parameters is reported in
\texttt{research/metric\_ansatz/} (8/8 PASS).
\end{corollary}

\begin{corollary}[Luminal gravitational waves]\label{cor:cGW}
Gravitational waves propagate on the effective metric $g_{\mathrm{eff}}$
at the vacuum speed: $c_{\mathrm{GW}}=c_{0}$ identically. This is
consistent with GW170817 constraints
($|c_{\mathrm{GW}}-c_{0}|/c_{0}<10^{-15}$).
\end{corollary}

\subsection{Field-dependent constants}

Requiring invariance of the Planck length $\lP=\sqrt{G\hbar/c^{3}}$
under the emergent rescaling fixes the field dependence of the
fundamental constants:
\begin{equation}\label{eq:constants}
  c(\Phi)=\frac{c_{0}}{\sqrt{\psi}},
  \qquad
  G(\Phi)=\frac{G_{0}}{\psi},
  \qquad
  \hbar(\Phi)=\frac{\hbar_{0}}{\sqrt{\psi}},
\end{equation}
all consistent with the exponential metric~\eqref{eq:exp-metric}. $\lP$
is the sole invariant across scales.

% =============================================================================
\section{Three physical regimes}\label{sec:regimes}
% =============================================================================

The non-linear potential $U(\varphi)$ together with the source term in
Eq.~\eqref{eq:field-eq} generates three distinct interaction regimes from a
single field equation. Linearization around a single static point source
yields a radial profile $\phi(r)$ with the following properties:

\begin{center}\small
\begin{tabular}{@{}lll@{}}
\toprule
\textbf{Regime} & \textbf{Scale} & \textbf{Physical content}\\
\midrule
(i)\ long-range attraction   & $r\gg r_{\mathrm{rep}}$ & Newtonian gravity, $U\sim -GM/r$\\
(ii)\ intermediate repulsion & $r\sim r_{\mathrm{rep}}$ & positive spatial potential between sources\\
(iii)\ short-range well      & $r\lesssim r_{\mathrm{well}}$ & confinement ($U\to-\infty$ finite well)\\
\bottomrule
\end{tabular}
\end{center}

The existence of three regimes follows from the non-monotonicity of
$\phi(r)$: a single generation profile acquires at least one inflection
point from the interplay of $\DD_{\mathrm{kin}}[\varphi]$ with the cubic
and quartic terms of $U(\varphi)$. The explicit scales $r_{\mathrm{rep}}$
and $r_{\mathrm{well}}$ depend on the source strength $qM/\PhiZero$ and on
the coupling $\gamma$ through
\begin{equation}\label{eq:Lambda-E}
  \Lambda_{E}\equiv\gamma/12,
  \qquad
  3H_{0}^{2}\approx c_{0}^{2}\,\Lambda_{E}
  \quad(\text{de Sitter limit}),
\end{equation}
which identifies $\Lambda_{E}$ as the effective cosmological constant.

\begin{remark}[Gravity as regime (i)]
Newtonian gravity is recovered as the long-range limit of the full
non-linear equation~\eqref{eq:field-eq}: in the weak field
$\Phi\approx\PhiZero(1+2U/c_{0}^{2})$, the leading source term reproduces
Poisson's equation $\nabla^{2}U=4\pi G\rho$ with
$G=q\PhiZero c_{0}^{2}/(8\pi)$ (identifying $q$, $\PhiZero$ with
Newton's constant).
\end{remark}

% =============================================================================
\section{Status of theorems and open problems}\label{sec:status}
% =============================================================================

We distinguish three levels of status, as follows:
\textsf{[TH]} formally proven theorem;
\textsf{[NU]} numerically verified to stated precision;
\textsf{[HY]} conjecture or hypothesis, not yet proven.
The list below summarises the current (April 2026) status of TGP core.

\subsection{Proven results (Core)}

\begin{center}\small
\begin{longtable}{@{}clp{7.2cm}c@{}}
\toprule
\# & \textbf{Result} & \textbf{Content} & \textbf{Status}\\
\midrule
\endhead
1 & Uniqueness of $\DD_{\mathrm{kin}}$
  & Theorem~\ref{thm:alpha2}: $K(\varphi)=\varphi^{4}$, $\alpha=2$ unique under (C1)--(C3)
  & \textsf{TH}\\
2 & Ghost-freedom (three sectors)
  & Sec.~\ref{ssec:ghost}: $K_{\mathrm{sub}}>0$ $\forall\varphi>0$; no kinetic singularity
  & \textsf{TH}\\
3 & Full field equation
  & Theorem~\ref{thm:field-eq}: Eq.~\eqref{eq:field-eq} from $\delta S/\delta\Phi=0$
  & \textsf{TH}\\
4 & Exponential metric uniqueness
  & Theorem~\ref{thm:metric}: PPN consistency $\Rightarrow$ exponential form
  & \textsf{TH}\\
5 & PPN parameters exact
  & Corollary~\ref{cor:ppn}: $\gamma_{\mathrm{PPN}}=\beta_{\mathrm{PPN}}=1$ to all orders
  & \textsf{TH}\\
6 & Luminal gravitational waves
  & Corollary~\ref{cor:cGW}: $c_{\mathrm{GW}}=c_{0}$ identically (GW170817)
  & \textsf{TH}\\
7 & Invariance of $\lP$
  & Eq.~\eqref{eq:constants}: $\lP=\sqrt{G\hbar/c^{3}}=\mathrm{const}$ consistent with $g_{\mu\nu}^{\mathrm{eff}}$
  & \textsf{TH}\\
8 & Instability of $\Nzero$
  & Proposition~\ref{prop:instability}: from Ising thermodynamics of $H_{\Gamma}$
  & \textsf{TH}\\
9 & $\Zp$ substrate symmetry
  & Proposition~\ref{prop:Z2}: $H_{\Gamma}$ invariant under $\hat{s}_{i}\mapsto-\hat{s}_{i}$
  & \textsf{TH}\\
10 & Three regimes from $U(\varphi)$
  & Sec.~\ref{sec:regimes}: non-monotonic $\phi(r)$ from cubic+quartic $U$
  & \textsf{TH}\\
11 & Bianchi propagation of Friedmann constraint
  & Remark following Eq.~(\ref{eq:field-eq}): $\nabla_{\mu}G_{\mathrm{eff}}^{\mu\nu}=0$
  & \textsf{TH}\\
12 & Dark-energy de~Sitter limit
  & Eq.~\eqref{eq:Lambda-E}: $3H_{0}^{2}\approx c_{0}^{2}\Lambda_{E}$, $\Lambda_{E}=\gamma/12$
  & \textsf{TH}\\
13 & Matter-source covariance
  & Axiom~\ref{ax:source} implemented by $\mathcal{L}_{\mathrm{mat}}=-q\varphi\rho/\PhiZero$
  & \textsf{TH}\\
14 & Vacuum condition $\beta=\gamma$
  & From $U'(1)=0$; Eq.~\eqref{eq:U-potential}
  & \textsf{TH}\\
15 & MS kinetic positivity in FRW
  & $Q_{s}=\varphi_{\mathrm{bg}}^{4}>0$ for all cosmological backgrounds
  & \textsf{TH}\\
\bottomrule
\end{longtable}
\end{center}

\subsection{Numerically verified (to stated precision)}

\begin{center}\small
\begin{longtable}{@{}clp{7.2cm}c@{}}
\toprule
\# & \textbf{Result} & \textbf{Content} & \textbf{Status}\\
\midrule
\endhead
16 & FRG structural stability of $K(\rho)\!\propto\!\rho$
  & One-loop LPA' flow near the Wilson--Fisher fixed point in $d=3$,
    $Z_{2}$ class: the kinetic form $K(\rho)\!\propto\!\rho$ is
    preserved, $K_{\mathrm{IR}}/K_{\mathrm{UV}}=1.000$, $\eta^{*}=0.044$.
    This verifies kinetic stability at the continuum level; it does
    \emph{not} substitute for the derivation of $\alpha$ from $H_{\Gamma}$
    (still OP-6).
  & \textsf{NU}\\
17 & Cosmology: $\Delta\chi^{2}=+1.8$, $\Delta\mathrm{AIC}=+3.8$
  & vs.\ $\Lambda$CDM on Planck\,+\,DESI data, Friedmann~\eqref{eq:field-eq}
  & \textsf{NU}\\
18 & Solitonic $N$-body stability
  & Chaos suppression, Earnshaw violation; 10/10 tests pass
  & \textsf{NU}\\
19 & Metric derivation chain
  & Spatial $\to$ information $\to$ PPN; cross-checked with $\lP=\mathrm{const}$
  & \textsf{NU}\\
20 & Three-regimes existence
  & Numerical profile $\phi(r)$ with attraction, repulsion, confinement
  & \textsf{NU}\\
21 & Coarse-graining $H_{\Gamma}\!\to\!S_{\mathrm{TGP}}$ (numerical)
  & Block-spin Monte Carlo on a 1D substrate chain
    (\texttt{research/continuum\_limit/cg\_strong\_numerical.py}) exhibits
    block-averaging behaviour consistent with a non-trivial fixed point,
    but does not determine the kinetic exponent. Earlier synthesis code
    in the same folder contained a derivation gap for $\alpha$ and has
    been withdrawn; see \texttt{KNOWN\_ISSUES.md}. The rigorous
    continuum limit remains OP-6; Theorem~\ref{thm:alpha2} is unaffected,
    being an axiomatic classification under (C1)--(C3).
  & \textsf{OP}\\
\bottomrule
\end{longtable}
\end{center}

\subsection{Open problems}

The following are explicitly open: they define the research frontier of
TGP at the core level.

\begin{center}\small
\begin{longtable}{@{}clp{8.5cm}c@{}}
\toprule
\# & \textbf{Open problem} & \textbf{Statement} & \textbf{Status}\\
\midrule
\endhead
OP-1 & Uniqueness of $U(\varphi)$
  & Is the cubic+quartic form the unique self-interference potential
  compatible with $\Zp$ symmetry and the vacuum condition?
  & \textsf{HY}\\
OP-2 & Parametric reduction $\beta=\gamma$
  & Formal proof (beyond the vacuum condition) that $\beta=\gamma$ is
  protected at all scales.
  & \textsf{HY}\\
OP-3 & Identification $a_{\Gamma}=1/\PhiZero$
  & Whether the two fitted parameters coincide, reducing the free input
  count to one.
  & \textsf{HY}\\
OP-4 & Derivation of the quartic coefficient $\gamma$
  & An \textit{a priori} determination of $\gamma$ (equivalently
  $\Lambda_{E}$) from $H_{\Gamma}$ alone.
  & \textsf{HY}\\
OP-5 & Full non-linear $N$-body regime
  & Existence/uniqueness of solutions for strongly overlapping sources
  beyond perturbation theory.
  & \textsf{HY}\\
OP-6 & Substrate metric continuum limit (rigorous)
  & Rigorous Migdal--Kadanoff convergence of $H_{\Gamma}$ to the continuum
  action~\eqref{eq:S-TGP} (Banach contraction of the block-spin map;
  $\Gamma$-convergence of the energy functional). Numerical evidence
  reported in row~21; fully rigorous proof open.
  & \textsf{HY}\\
OP-7 & Tensor projection $\sigma_{ab}$
  & Derivation of the explicit tensor-field equation from the same
  $H_{\Gamma}$; GW polarisation content.
  & \textsf{HY}\\
OP-8 & Gauge-sector emergence
  & Emergence of U(1)$\times$SU(2)$\times$SU(3) from the substrate
  (addressed partially in the extended manuscript).
  & \textsf{HY}\\
OP-9 & Strong-field black-hole solutions
  & Complete catalogue of non-singular BH solutions; quasi-normal modes.
  & \textsf{HY}\\
OP-10 & Quantum theory of $\Phi$
  & Second quantisation of $\Phi$ beyond the MS linear regime.
  & \textsf{HY}\\
\bottomrule
\end{longtable}
\end{center}

% =============================================================================
\section{Applications and falsifiability}\label{sec:applications}
% =============================================================================

The minimal core of TGP defined above motivates several companion
studies in which the same substrate and exponential metric are
confronted with separate empirical domains. Each companion paper cites
this reference for its foundational content; the items below are
pointers, not self-contained derivations.

\begin{itemize}[nosep]
  \item \textbf{Superconductivity.} Using the substrate harmonic
    $a^{\ast}_{\mathrm{TGP}}=7.725\,a_{\mathrm{B}}\approx 4.088$\,\AA\
    together with the orbital-amplitude spectrum
    $\{A_{s},A_{sp},A_{d},A_{f}\}$, the companion paper \cite{TGP-SC}
    presents a single-formula fit of $T_{c}$ across cuprates,
    phonon-mediated metals, pressure-driven systems and lanthanide
    hydrides in terms of three universal TGP constants plus tabulated
    atomic inputs. Scope, systematics and excluded families are
    discussed there.
  \item \textbf{Lepton masses.} A single substrate parameter
    $g_{0}^{e}=0.86941$ reproduces $m_{\mu}/m_{e}$ and $m_{\tau}/m_{e}$
    to better than $10^{-3}$, and the Koide combination $K=2/3$ is
    recovered as an arithmetic identity of the Brannen parametrisation
    \cite{TGP-leptons}; the identification of $g_{0}^{e}$ with a
    substrate quantity is a working hypothesis of that companion paper.
  \item \textbf{Black-hole shadows.} In the strong-field regime the
    exponential metric~\eqref{eq:exp-metric} leads to a photon-ring
    radius reduced by $\sim 37\%$ relative to Schwarzschild, offering a
    concrete target for next-generation EHT measurements
    \cite{TGP-BH}.
  \item \textbf{$N$-body dynamics.} Numerical integration of
    Eq.~\eqref{eq:field-eq} in configurations dominated by the
    intermediate-repulsion regime (regime~ii) exhibits chaos suppression
    and a violation of Earnshaw's theorem; see \cite{TGP-nbody} for the
    test suite.
\end{itemize}

\paragraph{Minimal falsification targets.} The following predictions are
of a binary character and provide the most direct falsification tests of
the core theory:
\begin{enumerate}[(F1),nosep]
  \item $\alpha=2$ \emph{identically} (Theorem~\ref{thm:alpha2}):
    any observable inconsistent with $K\propto\varphi^{4}$ at the
    substrate level falsifies TGP.
  \item $\gamma_{\mathrm{PPN}}=\beta_{\mathrm{PPN}}=1$ \emph{to all
    orders} (Corollary~\ref{cor:ppn}): a measured deviation at
    $O(U^{2}/c_{0}^{4})$ beyond GR would falsify the exponential metric.
  \item $c_{\mathrm{GW}}=c_{0}$ \emph{identically}
    (Corollary~\ref{cor:cGW}): already supported by GW170817 at the
    $10^{-15}$ level.
\end{enumerate}

\paragraph{Closed-branch negative checks.} Three potential near-term
falsification channels have been analysed in dedicated companion
studies and found to lie \emph{below} current experimental precision,
closing them as viable discovery avenues without contradicting the
theory:
\begin{itemize}[nosep]
  \item \emph{Muon $(g\!-\!2)$:} the TGP-substrate correction to the
    anomalous magnetic moment is $\Delta a_{\mu}\lesssim 10^{-12}$,
    three orders of magnitude below the current experimental
    uncertainty; see \texttt{research/muon\_g\_minus\_2/}.
  \item \emph{Wiedemann--Franz in bad metals:} the substrate-diffusion
    correction to $L/L_{0}$ is $\mathcal{O}(10^{-4})$ at $300$\,K, three
    to four orders of magnitude smaller than observed deviations in
    cuprates and heavy-fermion compounds, which are therefore \emph{not}
    of TGP origin; see \texttt{research/thermal\_transport\_molecular/}.
  \item \emph{Casimir force in MOFs:} the TGP correction is at the
    $10^{-4}$--$10^{-3}$ level, below current torsion-balance precision;
    see \texttt{research/casimir\_mof/}.
\end{itemize}
These null results are stated not as corroborations but as explicit
constraints on where a TGP-specific signature could plausibly emerge
with present instrumentation.

% =============================================================================
\section{Conclusion}\label{sec:conclusion}
% =============================================================================

We have presented a minimal, self-contained core of the Theory of Generated
Space: four axioms, one scalar field with a geometrically fixed kinetic
operator ($\alpha=2$), a single non-linear potential with two coefficients
related by a vacuum condition, and an algebraically determined effective
metric reproducing general relativity exactly at the post-Newtonian level.
Fifteen core results are formally proven, six more are numerically
verified, and ten explicit open problems are stated so that the theory is
both well-defined and falsifiable on a clear basis.

This core is intended to serve as a common reference for applications to
superconductivity, particle masses, black-hole phenomenology and cosmology,
developed in companion papers.

\paragraph{Data and reproducibility.} All numerical scripts underlying the
results of this paper are available in the accompanying source repository
\cite{TGP-repo}. The present paper deposits a fixed, versioned snapshot on
Zenodo (DOI assigned at publication).

\paragraph{Acknowledgements.} The author thanks the community of
open-source physics for the tools that made this work possible.

% =============================================================================
\begin{thebibliography}{99}
% =============================================================================

\bibitem{TGP-repo}
M.~Serafin, \emph{Theory of Generated Space --- Core paper source repository},
\url{https://github.com/Stefan13610/tgp-core-paper} (2026).

\bibitem{TGP-SC}
M.~Serafin,
\emph{Unified superconductivity from the TGP substrate},
companion paper, Zenodo (2026).

\bibitem{TGP-leptons}
M.~Serafin,
\emph{Lepton masses and the Koide relation from a single TGP parameter},
companion paper (2026).

\bibitem{TGP-BH}
M.~Serafin,
\emph{Black-hole shadows in the TGP exponential metric},
companion paper (2026).

\bibitem{TGP-nbody}
M.~Serafin,
\emph{N-body dynamics in TGP: chaos suppression and Earnshaw violation},
companion paper (2026).

\bibitem{Will2014}
C.~M.~Will,
\emph{The confrontation between general relativity and experiment},
Living Rev.\ Relativity \textbf{17}, 4 (2014).

\bibitem{BertottiCassini}
B.~Bertotti, L.~Iess, and P.~Tortora,
\emph{A test of general relativity using radio links with the Cassini spacecraft},
Nature \textbf{425}, 374 (2003).

\bibitem{WilliamsLLR}
J.~G.~Williams, S.~G.~Turyshev, and D.~H.~Boggs,
\emph{Progress in lunar laser ranging tests of relativistic gravity},
Phys.\ Rev.\ Lett.\ \textbf{93}, 261101 (2004).

\bibitem{Planck2018}
Planck Collaboration,
\emph{Planck 2018 results.\ VI.\ Cosmological parameters},
A\&A \textbf{641}, A6 (2020).

\bibitem{GW170817-speed}
LIGO/Virgo Collaboration,
\emph{Gravitational waves and gamma-rays from a binary neutron star merger:
GW170817 and GRB 170817A}, ApJL \textbf{848}, L13 (2017).

\end{thebibliography}

\end{document}
